% GENERATED FILE. Edit canonical lessons, then run npm run generate:books.
% canonical-source-hash: fnv1a64:27f10471054d7aa9
% canonical-lessons: HI-C167-kiraya, HI-C167-tarikh, HI-C167-nambar, HI-C167-manzil, HI-C167-rasid, HI-R167-first-pass-notices-and-messages, HI-R167-second-pass-rent-dates-and-rooms

\chapter{Rent, a Date, a Number, a Floor, a Receipt}
\label{ch:hi-rent-a-date-a-number-a-floor-a-receipt}
\hlchaptermodality{167}

\begin{chapteropening}
\textbf{By the end of this chapter:} \emph{I can read the rent, a date, a room number, a floor and a receipt in a short notice.}

Complete the last lesson of chapter 167: I can read the rent, a date, a room number, a floor and a receipt in a short notice.
\end{chapteropening}

\section[\texorpdfstring{\hi{किराया} (kirāyā)}{किराया (kirāyā)}]{\hi{किराया} (kirāyā) — rent, a fare}

\label{lesson:HI-C167-kiraya}



\begin{quote}
Before the new one: say the Hindi for an advertisement, then the Hindi for a repair.
\end{quote}

\subsection*{You'll want to know: \hi{किराया}}
\textbf{\hi{किराया}} — \emph{kirāyā} — \textquotedblleft{}rent, a fare\textquotedblright{}.

\emph{Read it:} the advert

The advert says \textbf{\hi{कमरा} \hi{ख़ाली} \hi{है।} \hi{किराया}: \hi{पाँच} \hi{सौ} \hi{रुपये।}} (\emph{kamrā khālī hai. kirāyā: pā̃c sau rupaye.}) The rent is five hundred rupees.

\begin{grammarlens}[title={What you've built}]
One more word for this chapter.
\end{grammarlens}

\subsection*{Your turn}
\begin{itemize}
\raggedright
  \item \emph{Say it:} \emph{kirāyā}
  \item \emph{Say it:} \emph{kirāyā}, once more
  \item \emph{Say it:} say \emph{kirāyā}
  \item \emph{Recall:} say the Hindi for an advertisement, then the Hindi for a repair, then say \emph{kirāyā} again
\end{itemize}

\begin{tcolorbox}[breakable,colback=teal!4,colframe=teal!35!black,title={Before you move on}]
What does \hi{किराया} mean? (\textquotedblleft{}Rent, a fare\textquotedblright{}.) Say it once more.
\end{tcolorbox}
\section[\texorpdfstring{\hi{तारीख़} (tārīkh)}{तारीख़ (tārīkh)}]{\hi{तारीख़} (tārīkh) — a date}

\label{lesson:HI-C167-tarikh}



\begin{quote}
Before the new one: say the Hindi for a repair, then the Hindi for rent.
\end{quote}

\subsection*{You'll want to know: \hi{तारीख़}}
\textbf{\hi{तारीख़}} — \emph{tārīkh} — \textquotedblleft{}a date\textquotedblright{}.

\emph{Read it:} the note

The note says \textbf{\hi{तारीख़} \hi{पाँच}, \hi{सुबह} \hi{नौ} \hi{बजे।}} (\emph{tārīkh pā̃c, subah nau baje.}) The fifth, at nine in the morning.

\begin{grammarlens}[title={What you've built}]
One more word for this chapter.
\end{grammarlens}

\subsection*{Your turn}
\begin{itemize}
\raggedright
  \item \emph{Say it:} \emph{tārīkh}
  \item \emph{Say it:} \emph{tārīkh}, once more
  \item \emph{Say it:} say \emph{tārīkh}
  \item \emph{Recall:} say the Hindi for a repair, then the Hindi for rent, then say \emph{tārīkh} again
\end{itemize}

\begin{tcolorbox}[breakable,colback=teal!4,colframe=teal!35!black,title={Before you move on}]
What does \hi{तारीख़} mean? (\textquotedblleft{}A date\textquotedblright{}.) Say it once more.
\end{tcolorbox}
\section[\texorpdfstring{\hi{नंबर} (nambar)}{नंबर (nambar)}]{\hi{नंबर} (nambar) — a number}

\label{lesson:HI-C167-nambar}



\begin{quote}
Before the new one: say the Hindi for rent, then the Hindi for a date.
\end{quote}

\subsection*{You'll want to know: \hi{नंबर}}
\textbf{\hi{नंबर}} — \emph{nambar} — \textquotedblleft{}a number\textquotedblright{}.

\emph{Read it:} the sign

The sign says \textbf{\hi{कमरा} \hi{नंबर} \hi{बीस।}} (\emph{kamrā nambar bīs.}) Room twenty.

\begin{grammarlens}[title={What you've built}]
One more word for this chapter.
\end{grammarlens}

\subsection*{Your turn}
\begin{itemize}
\raggedright
  \item \emph{Say it:} \emph{nambar}
  \item \emph{Say it:} \emph{nambar}, once more
  \item \emph{Say it:} say \emph{nambar}
  \item \emph{Recall:} say the Hindi for rent, then the Hindi for a date, then say \emph{nambar} again
\end{itemize}

\begin{tcolorbox}[breakable,colback=teal!4,colframe=teal!35!black,title={Before you move on}]
What does \hi{नंबर} mean? (\textquotedblleft{}A number\textquotedblright{}.) Say it once more.
\end{tcolorbox}
\section[\texorpdfstring{\hi{मंज़िल} (manzil)}{मंज़िल (manzil)}]{\hi{मंज़िल} (manzil) — a floor, a storey}

\label{lesson:HI-C167-manzil}



\begin{quote}
Before the new one: say the Hindi for a date, then the Hindi for a number.
\end{quote}

\subsection*{You'll want to know: \hi{मंज़िल}}
\textbf{\hi{मंज़िल}} — \emph{manzil} — \textquotedblleft{}a floor, a storey\textquotedblright{}.

\emph{Read it:} the sign

The sign says \textbf{\hi{कमरा} \hi{नंबर} \hi{बीस}, \hi{दूसरी} \hi{मंज़िल।}} (\emph{kamrā nambar bīs, dūsrī manzil.}) Room twenty is on the second floor: go up.

\begin{grammarlens}[title={What you've built}]
One more word for this chapter.
\end{grammarlens}

\subsection*{Your turn}
\begin{itemize}
\raggedright
  \item \emph{Say it:} \emph{manzil}
  \item \emph{Say it:} \emph{manzil}, once more
  \item \emph{Say it:} say \emph{manzil}
  \item \emph{Recall:} say the Hindi for a date, then the Hindi for a number, then say \emph{manzil} again
\end{itemize}

\begin{tcolorbox}[breakable,colback=teal!4,colframe=teal!35!black,title={Before you move on}]
What does \hi{मंज़िल} mean? (\textquotedblleft{}A floor, a storey\textquotedblright{}.) Say it once more.
\end{tcolorbox}
\section[\texorpdfstring{\hi{रसीद} (rasīd)}{रसीद (rasīd)}]{\hi{रसीद} (rasīd) — a receipt}

\label{lesson:HI-C167-rasid}



\begin{quote}
Before the new one: say the Hindi for a number, then the Hindi for a floor.
\end{quote}

\subsection*{You'll want to know: \hi{रसीद}}
\textbf{\hi{रसीद}} — \emph{rasīd} — \textquotedblleft{}a receipt\textquotedblright{}.

\emph{Read it:} the receipt

The receipt says \textbf{\hi{रसीद}: \hi{किराया} \hi{पाँच} \hi{सौ} \hi{रुपये}, \hi{तारीख़} \hi{पाँच।}} (\emph{rasīd: kirāyā pā̃c sau rupaye, tārīkh pā̃c.}) Five hundred rupees of rent, paid on the fifth: keep it.

\begin{grammarlens}[title={What you've built}]
One more word for this chapter.
\end{grammarlens}

\subsection*{Your turn}
\begin{itemize}
\raggedright
  \item \emph{Say it:} \emph{rasīd}
  \item \emph{Say it:} \emph{rasīd}, once more
  \item \emph{Say it:} say \emph{rasīd}
  \item \emph{Recall:} say the Hindi for a number, then the Hindi for a floor, then say \emph{rasīd} again
\end{itemize}

\begin{tcolorbox}[breakable,colback=teal!4,colframe=teal!35!black,title={Before you move on}]
What does \hi{रसीद} mean? (\textquotedblleft{}A receipt\textquotedblright{}.) Say it once more.
\end{tcolorbox}
\section[(dialogue)]{Notices and messages — a first pass}

\label{lesson:HI-R167-first-pass-notices-and-messages}



\begin{quote}
Say the words for a floor and a receipt.
\end{quote}

\subsection*{Your turn}
\begin{itemize}
\raggedright
  \item \emph{Say it:} a notice, a message, time, an advertisement, a repair
\end{itemize}

\begin{tcolorbox}[breakable,colback=teal!4,colframe=teal!35!black,title={Before you move on}]
A notice? (\textbf{\hi{सूचना}}, \emph{sūcnā}.) A message? (\textbf{\hi{संदेश}}, \emph{sandeś}.) Time? (\textbf{\hi{समय}}, \emph{samay}.) An advertisement? (\textbf{\hi{विज्ञापन}}, \emph{vijñāpan}.) A repair? (\textbf{\hi{मरम्मत}}, \emph{marammat}.)
\end{tcolorbox}
\section[(dialogue)]{Rent, dates and rooms — a second pass}

\label{lesson:HI-R167-second-pass-rent-dates-and-rooms}



\begin{quote}
Say the words for rent and a date.
\end{quote}

\subsection*{Your turn}
\begin{itemize}
\raggedright
  \item \emph{Say it:} rent, a date, a number, a floor, a receipt
\end{itemize}

\begin{tcolorbox}[breakable,colback=teal!4,colframe=teal!35!black,title={Before you move on}]
Rent? (\textbf{\hi{किराया}}, \emph{kirāyā}.) A date? (\textbf{\hi{तारीख़}}, \emph{tārīkh}.) A number? (\textbf{\hi{नंबर}}, \emph{nambar}.) A floor? (\textbf{\hi{मंज़िल}}, \emph{manzil}.) A receipt? (\textbf{\hi{रसीद}}, \emph{rasīd}.)
\end{tcolorbox}
